\documentclass[11pt,a4paper]{article}
\usepackage[utf8]{inputenc}
\usepackage[spanish,es-tabla]{babel}
\usepackage{amsmath,amsfonts,amssymb}
\usepackage{geometry}
\usepackage{graphicx}
\usepackage{xcolor}

\usepackage{enumitem}
\usepackage{tikz}

\geometry{margin=2.5cm}

\title{\textbf{Soluci\'on Examen 2025-2026\\TAF STAR - UEC B\\14/10/2025}}
\author{}
\date{}

\begin{document}

\maketitle

\section*{Preguntas con respuestas cortas (6 puntos)}

\subsection*{Pregunta 1: Receptor \'optimo para modulaci\'on lineal (2 pts)}

\textbf{Enunciado:} Consideramos un receptor \'optimo para una modulaci\'on lineal sobre frecuencia portadora transmitida sobre un canal gaussiano. La respuesta impulsional del filtro pasa-bajo de emisi\'on es $h(t)$. Se denota por $T_s$ el periodo s\'imbolo. En el receptor, el $n$-\'esimo instante de muestreo es igual a $t_0 + nT_s$.

\textbf{a)} Sea $g(t)$ la respuesta impulsional del filtro de recepci\'on cuya salida se muestrea. Dar la expresi\'on de $g(t)$ en funci\'on de las hip\'otesis y notaciones.

\textbf{b)} ¿Cu\'ales son los dos roles principales de este filtro?

\vspace{0.3cm}
\textbf{Soluci\'on:}

\textbf{a) Expresi\'on de $g(t)$:}

Para un receptor \'optimo de m\'axima verosimilitud (ML) en un canal gaussiano, el filtro de recepci\'on debe ser el \textbf{filtro adaptado} al filtro de emisi\'on:

\[
\boxed{g(t) = h(t_0 - t)}
\]

Donde:
\begin{itemize}
    \item $h(t)$ es la respuesta impulsional del filtro de emisi\'on
    \item $t_0$ es un retardo que garantiza la causalidad del filtro
    \item El filtro $g(t)$ es una versi\'on ''volteada en el tiempo'' y desplazada de $h(t)$
\end{itemize}

\textbf{Justificaci\'on matem\'atica:}

El filtro adaptado maximiza el SNR en el instante de muestreo. Si la se\~nal recibida es $r(t) = s(t) + n(t)$ con $s(t) = \sum_k a_k h(t-kT_s)\cos(2\pi\nu_0 t + \theta_0)$, entonces el filtro $g(t) = h(t_0 - t)$ es \'optimo.

\textbf{b) Dos roles principales del filtro:}

\textbf{Rol 1: Maximizaci\'on del SNR}

El filtro adaptado maximiza la relaci\'on se\~nal a ruido en el instante de muestreo $t_0 + nT_s$.

\textit{Demostraci\'on:} Por la desigualdad de Cauchy-Schwarz, el SNR m\'aximo se obtiene cuando el filtro es proporcional a $h(t_0 - t)$.

\textbf{Rol 2: Control de la interferencia entre s\'imbolos (ISI)}

La cascada del filtro de emisi\'on y el filtro de recepci\'on debe satisfacer el criterio de Nyquist para eliminar la ISI:
\[
(h * g)(t_0 + kT_s) = \delta_{k,0}
\]

Esto equivale al criterio TSWOC:
\[
\int_{\mathbb{R}} h(t) h(t - kT_s)\, dt = \delta_{k,0}
\]

\textbf{Consecuencias pr\'acticas:}

\begin{itemize}
    \item \textbf{Maximizaci\'on del SNR:} Mejora el rendimiento del sistema (menor probabilidad de error)
    
    \item \textbf{Ausencia de ISI:} Permite detecci\'on s\'imbolo por s\'imbolo sin necesidad de ecualizadores complejos
    
    \item \textbf{Detecci\'on \'optima:} Combinando ambas propiedades, se logra la m\'inima probabilidad de error posible en un canal AWGN
\end{itemize}

\textbf{Implementaci\'on pr\'actica:}

En la pr\'actica, se usan filtros raised cosine:
\begin{itemize}
    \item Emisi\'on: root raised cosine $h_{RRC}(t)$
    \item Recepci\'on: root raised cosine $g_{RRC}(t) = h_{RRC}(t_0 - t)$
    \item Cascada: $(h_{RRC} * g_{RRC})(t)$ = raised cosine que satisface Nyquist
\end{itemize}



\subsection*{Pregunta 2: Eficiencia espectral de 256-QAM (2 pts)}

\textbf{Enunciado:} Consideremos una modulaci\'on 256-QAM. La banda ocupada se elige igual a la banda m\'inima necesaria para transmitir sin interferencia entre s\'imbolos (ISI). Calcular la eficiencia espectral.

\vspace{0.3cm}
\textbf{Soluci\'on:}

\textbf{Datos:}
\begin{itemize}
    \item Modulaci\'on: 256-QAM
    \item Orden: $M = 256 = 2^8$
    \item Banda ocupada: m\'inima para transmisi\'on sin ISI
\end{itemize}

\textbf{Desarrollo:}

\textbf{Paso 1: Banda m\'inima para transmisi\'on sin ISI}

La banda m\'inima (criterio de Nyquist) para una modulaci\'on sobre frecuencia portadora es:
\[
B_{occ,min} = \frac{1}{T_s} = D_s
\]

Donde:
\begin{itemize}
    \item $T_s$ es el periodo s\'imbolo
    \item $D_s$ es la rapidez de s\'imbolo (symbol rate)
\end{itemize}

Esto corresponde a un filtro rectangular ideal (roll-off $\beta = 0$).

\textbf{Paso 2: D\'ebito binario}

El d\'ebito binario se relaciona con la rapidez de s\'imbolo mediante:
\[
D_b = D_s \times \log_2(M)
\]

Para 256-QAM:
\[
D_b = D_s \times \log_2(256) = D_s \times 8
\]

\textbf{Paso 3: Eficiencia espectral}

La eficiencia espectral se define como:
\[
\eta = \frac{D_b}{B_{occ}} \text{ [bits/s/Hz]}
\]

Sustituyendo $B_{occ} = D_s$ y $D_b = 8 D_s$:

\[
\eta = \frac{8 D_s}{D_s} = \boxed{8 \text{ bits/s/Hz}}
\]

\textbf{Interpretaci\'on:}

\begin{itemize}
    \item \textbf{Valor m\'aximo te\'orico:} Para la banda m\'inima de Nyquist, la eficiencia espectral es:
    \[
    \eta_{max} = \log_2(M) = 8 \text{ bits/s/Hz}
    \]
    
    \item \textbf{Comparaci\'on con otras modulaciones:}
    \begin{center}
    \begin{tabular}{|c|c|c|}
    \hline
    \textbf{Modulaci\'on} & \textbf{$\log_2(M)$} & \textbf{$\eta$ [bits/s/Hz]} \\
    \hline
    BPSK & 1 & 1 \\
    QPSK & 2 & 2 \\
    16-QAM & 4 & 4 \\
    64-QAM & 6 & 6 \\
    256-QAM & 8 & 8 \\
    1024-QAM & 10 & 10 \\
    \hline
    \end{tabular}
    \end{center}
    
    \item \textbf{Pr\'actica vs. teor\'ia:}
    
    En sistemas reales, se usa roll-off $\beta > 0$ (t\'ipicamente $\beta = 0.22$ en LTE, $\beta = 0.25$ en 5G):
    \[
    B_{occ,real} = (1 + \beta) D_s
    \]
    
    La eficiencia espectral real es:
    \[
    \eta_{real} = \frac{\log_2(M)}{1 + \beta}
    \]
    
    Para 256-QAM con $\beta = 0.25$:
    \[
    \eta_{real} = \frac{8}{1.25} = 6.4 \text{ bits/s/Hz}
    \]
    
    \item \textbf{L\'imite de Shannon:}
    
    La capacidad de Shannon para un canal AWGN es:
    \[
    C = W \log_2\left(1 + \frac{S}{N}\right)
    \]
    
    La eficiencia espectral est\'a limitada por:
    \[
    \eta \leq \log_2\left(1 + \text{SNR}\right)
    \]
    
    Para $\eta = 8$ bits/s/Hz, se requiere:
    \[
    \text{SNR} \geq 2^8 - 1 = 255 \approx 24 \text{ dB}
    \]
\end{itemize}

\textbf{Aplicaciones pr\'acticas:}

256-QAM se usa en:
\begin{itemize}
    \item LTE-Advanced (downlink en buenas condiciones)
    \item 5G NR (con SNR alto)
    \item Cable modem DOCSIS 3.1
    \item Wi-Fi 802.11ac/ax (en enlaces con SNR muy alto)
\end{itemize}

\textbf{Conclusi\'on:} La eficiencia espectral de 256-QAM con banda m\'inima de Nyquist es $\boxed{8 \text{ bits/s/Hz}}$, que es la m\'axima te\'orica para esta modulaci\'on.



\newpage
\section*{Transmisi\'on 16-ASK sobre canal AWGN (14 puntos)}

\textbf{Datos del problema:}
\begin{itemize}
    \item Elementos binarios independientes e id\'enticamente distribuidos: $\Pr(m_i = 0) = \Pr(m_i = 1)$
    \item Modulaci\'on: 16-ASK
\end{itemize}

\textbf{Se\~nal modulada:}
\[
x(t) = \sum_{k=0}^{L-1} a_k h(t-kT_s)\cos(2\pi\nu_0 t + \theta_0)
\]

\textbf{Funci\'on de transferencia del filtro:}
\[
H(\nu) = \begin{cases}
\sqrt{T_s} & |\nu| \leq \frac{1}{4T_s} \\
\sqrt{2T_s\sqrt{\frac{3}{4} - T_s|\nu|}} & \frac{1}{4T_s} \leq |\nu| \leq \frac{3}{4T_s} \\
0 & \text{en otro caso}
\end{cases}
\]

\textbf{Par\'ametros:}
\begin{itemize}
    \item $T_s = 10^{-6}$ s
    \item $\nu_0 = 10$ MHz = $10^7$ Hz
\end{itemize}

\textbf{Canal:} Gaussiano. Se\~nal recibida: $r(t) = x(t) + w(t)$ con $w(t)$ ruido AWGN de media cero y densidad espectral de potencia $\frac{N_0}{2}$.

\textbf{Receptor:}

\begin{center}
\begin{tikzpicture}[node distance=2.5cm]
\node (input) {$r(t)$};
\node[draw, right of=input, node distance=2cm] (mult) {$\times \cos(2\pi\nu_0 t + \theta_1)$};
\node[draw, right of=mult] (filter) {$g(t)$};
\node[draw, right of=filter, node distance=2cm] (sampler) {$t_0+nT_s$};
\node[right of=sampler, node distance=2cm] (output) {$z_n$};

\draw[->] (input) -- (mult);
\draw[->] (mult) -- (filter);
\draw[->] (filter) -- (sampler);
\draw[->] (sampler) -- (output);
\end{tikzpicture}
\end{center}

El valor $\theta_1$ puede diferir de $\theta_0$ con $\theta_1 - \theta_0 \in [0, \frac{\pi}{2}]$.

\subsection*{a) Definir el alfabeto elemental 16-ASK (1 pt)}

\textbf{Enunciado:} Definir el alfabeto elemental 16-ASK.

\vspace{0.3cm}
\textbf{Soluci\'on:}

El alfabeto 16-ASK consiste en 16 niveles de amplitud equiespaciados y sim\'etricos:

\[
\boxed{\mathcal{A} = \{-15, -13, -11, -9, -7, -5, -3, -1, +1, +3, +5, +7, +9, +11, +13, +15\}}
\]

\textbf{Propiedades:}
\begin{itemize}
    \item \textbf{N\'umero de s\'imbolos:} $M = 16$
    \item \textbf{Bits por s\'imbolo:} $\log_2(16) = 4$ bits
    \item \textbf{Separaci\'on entre s\'imbolos:} $\Delta = 2$
    \item \textbf{Media:} $\mathbb{E}[a_k] = 0$ (alfabeto sim\'etrico)
    \item \textbf{Energ\'ia media:}
    \begin{align*}
    \mathbb{E}[a_k^2] &= \frac{1}{16}(15^2 + 13^2 + 11^2 + 9^2 + 7^2 + 5^2 + 3^2 + 1^2) \times 2 \\
    &= \frac{1}{8}(225 + 169 + 121 + 81 + 49 + 25 + 9 + 1) \\
    &= \frac{680}{8} = 85
    \end{align*}
    \item \textbf{Distancia m\'inima:} $d_{min} = 2$
    \item \textbf{Rango:} $[-15, +15]$
\end{itemize}

\textbf{F\'ormula general para M-ASK:}

Para un alfabeto M-ASK con $M = 2^m$ s\'imbolos:
\[
\mathcal{A} = \{-(M-1), -(M-3), \ldots, -1, +1, \ldots, (M-3), (M-1)\}
\]

Energ\'ia media:
\[
\mathbb{E}[a^2] = \frac{(M^2 - 1)}{3}
\]

Para $M = 16$: $\mathbb{E}[a^2] = \frac{256 - 1}{3} = \frac{255}{3} = 85$ ✓



\subsection*{b) Proponer un Gray mapping (2 pts)}

\textbf{Enunciado:} Proponer un Gray mapping para la conversi\'on binario a s\'imbolo.

\vspace{0.3cm}
\textbf{Soluci\'on:}

\textbf{Gray mapping propuesto:}

\begin{center}
\begin{tabular}{|c|c|c|c|c|c|}
\hline
\textbf{$a_k$} & \textbf{$m_{4k}$} & \textbf{$m_{4k+1}$} & \textbf{$m_{4k+2}$} & \textbf{$m_{4k+3}$} & \textbf{C\'odigo} \\
\hline
$-15$ & 0 & 0 & 0 & 0 & 0000 \\
$-13$ & 0 & 0 & 0 & 1 & 0001 \\
$-11$ & 0 & 0 & 1 & 1 & 0011 \\
$-9$ & 0 & 0 & 1 & 0 & 0010 \\
$-7$ & 0 & 1 & 1 & 0 & 0110 \\
$-5$ & 0 & 1 & 1 & 1 & 0111 \\
$-3$ & 0 & 1 & 0 & 1 & 0101 \\
$-1$ & 0 & 1 & 0 & 0 & 0100 \\
$+1$ & 1 & 1 & 0 & 0 & 1100 \\
$+3$ & 1 & 1 & 0 & 1 & 1101 \\
$+5$ & 1 & 1 & 1 & 1 & 1111 \\
$+7$ & 1 & 1 & 1 & 0 & 1110 \\
$+9$ & 1 & 0 & 1 & 0 & 1010 \\
$+11$ & 1 & 0 & 1 & 1 & 1011 \\
$+13$ & 1 & 0 & 0 & 1 & 1001 \\
$+15$ & 1 & 0 & 0 & 0 & 1000 \\
\hline
\end{tabular}
\end{center}

\textbf{Verificaci\'on del criterio Gray:}

Cada par de s\'imbolos adyacentes debe diferir en exactamente 1 bit:
\begin{itemize}
    \item $-15 \rightarrow -13$: $0000 \rightarrow 0001$ (difieren en bit 3) ✓
    \item $-13 \rightarrow -11$: $0001 \rightarrow 0011$ (difieren en bit 2) ✓
    \item $-11 \rightarrow -9$: $0011 \rightarrow 0010$ (difieren en bit 3) ✓
    \item $-9 \rightarrow -7$: $0010 \rightarrow 0110$ (difieren en bit 1) ✓
    \item Y as\'i sucesivamente...
\end{itemize}

\textbf{Estructura del c\'odigo:}

\begin{itemize}
    \item \textbf{Bit $m_{4k}$ (MSB):} Indica el signo
    \begin{itemize}
        \item $0$: s\'imbolos negativos ($a_k < 0$)
        \item $1$: s\'imbolos positivos ($a_k > 0$)
    \end{itemize}
    
    \item \textbf{Bits $m_{4k+1}$, $m_{4k+2}$, $m_{4k+3}$:} Codifican la magnitud usando Gray code de 3 bits
\end{itemize}

\textbf{Ventaja:} Minimiza la tasa de error binaria cuando los errores ocurren preferentemente entre s\'imbolos adyacentes (caso t\'ipico en AWGN con SNR moderada a alta).



\subsection*{c) Calcular y trazar $a(t)$ (1 pt)}

\textbf{Enunciado:} Definimos la se\~nal $a(t) = \sum_{k=0}^{L-1} a_k \delta(t - kT_s)$. Debemos modular el tren binario $m = [0000\ 1010\ 1100\ 0101]$. Calcular y trazar $a(t)$.

\vspace{0.3cm}
\textbf{Soluci\'on:}

\textbf{Conversi\'on binaria a s\'imbolos:}

Dividimos el tren binario en grupos de 4 bits:
\begin{itemize}
    \item Bits $[0000]$: seg\'un la tabla, corresponde a $a_0 = -15$
    \item Bits $[1010]$: seg\'un la tabla, corresponde a $a_1 = +9$
    \item Bits $[1100]$: seg\'un la tabla, corresponde a $a_2 = +1$
    \item Bits $[0101]$: seg\'un la tabla, corresponde a $a_3 = -3$
\end{itemize}

\textbf{Expresi\'on de $a(t)$:}

\[
\boxed{a(t) = -15\delta(t) + 9\delta(t - T_s) + 1\delta(t - 2T_s) - 3\delta(t - 3T_s)}
\]

\textbf{Representaci\'on gr\'afica:}

\begin{center}
\begin{tikzpicture}[scale=0.8]
% Eje temporal
\draw[->] (-0.5,0) -- (5,0) node[right] {$t$};
\draw[->] (0,-2) -- (0,2) node[above] {$a(t)$};

% Marcas temporales
\foreach \x/\label in {0/0, 1/T_s, 2/2T_s, 3/3T_s}
{
    \node[below] at (\x,-0.1) {$\label$};
    \draw (\x,-0.05) -- (\x,0.05);
}

% Impulsos (flechas hacia arriba/abajo)
% a_0 = -15 en t=0
\draw[->, very thick, red] (0,0) -- (0,-1.5) node[left, pos=0.5] {$-15$};
\fill[red] (0,0) circle (2pt);

% a_1 = +9 en t=T_s
\draw[->, very thick, blue] (1,0) -- (1,0.9) node[right, pos=0.5] {$+9$};
\fill[blue] (1,0) circle (2pt);

% a_2 = +1 en t=2T_s
\draw[->, very thick, green!70!black] (2,0) -- (2,0.1) node[right, pos=0.5] {$+1$};
\fill[green!70!black] (2,0) circle (2pt);

% a_3 = -3 en t=3T_s
\draw[->, very thick, orange] (3,0) -- (3,-0.3) node[left, pos=0.5] {$-3$};
\fill[orange] (3,0) circle (2pt);

% C\'odigos binarios
\node[red, above] at (0,0.3) {0000};
\node[blue, above] at (1,1.2) {1010};
\node[green!70!black, above] at (2,0.4) {1100};
\node[orange, below] at (3,-0.6) {0101};

\end{tikzpicture}
\end{center}

\textbf{Interpretaci\'on:}

La se\~nal $a(t)$ consiste en 4 impulsos de Dirac (deltas) en los instantes $t = 0, T_s, 2T_s, 3T_s$ con pesos (amplitudes) $-15, +9, +1, -3$ respectivamente.

\textbf{Nota:} En la pr\'actica, la se\~nal transmitida es:
\[
x(t) = \sum_{k=0}^{3} a_k h(t-kT_s)\cos(2\pi\nu_0 t + \theta_0)
\]

donde cada impulso $a_k\delta(t-kT_s)$ es conformado por el filtro $h(t)$ y modulado en la portadora.



\subsection*{d) Demostrar que el filtro satisface TSWOC (2 pts)}

\textbf{Enunciado:} Demostrar que el filtro de transmisi\'on (definido por su funci\'on de transferencia $H(\nu)$) satisface el criterio TSWOC. Se espera una expresi\'on cuidadosa del criterio y la demostraci\'on puede basarse en una representaci\'on gr\'afica relevante.

\vspace{0.3cm}
\textbf{Soluci\'on:}

\textbf{Criterio TSWOC en el dominio frecuencial:}

\[
\boxed{\sum_k |H(\nu - k/T_s)|^2 = T_s, \quad \forall \nu \in \mathbb{R}}
\]

O equivalentemente en el dominio temporal:
\[
\int_{\mathbb{R}} h(t) h(t - kT_s)\, dt = \delta_{k,0}
\]

\textbf{Demostraci\'on gr\'afica:}

Necesitamos verificar que la suma $P(\nu) = \sum_k |H(\nu - k/T_s)|^2$ es constante e igual a $T_s$.

\textbf{An\'alisis por intervalos:}

Con $T_s = 10^{-6}$ s, tenemos $\frac{1}{T_s} = 10^6$ Hz.

El soporte de $H(\nu)$ es $\left[-\frac{3}{4T_s}, \frac{3}{4T_s}\right] = [-0.75 \times 10^6, 0.75 \times 10^6]$ Hz.

\textbf{Intervalo 1:} $\nu \in \left[\frac{p}{T_s} - \frac{1}{4T_s}, \frac{p}{T_s} + \frac{1}{4T_s}\right]$

Solo $|H(\nu - p/T_s)|^2$ es no nulo:
\[
P(\nu) = |H(\nu - p/T_s)|^2 = T_s \quad ✓
\]

\textbf{Intervalo 2:} $\nu \in \left[\frac{p}{T_s} + \frac{1}{4T_s}, \frac{p}{T_s} + \frac{3}{4T_s}\right]$

Son no nulos $|H(\nu - p/T_s)|^2$ y $|H(\nu - (p+1)/T_s)|^2$:

\begin{align*}
P(\nu) &= |H(\nu - p/T_s)|^2 + |H(\nu - (p+1)/T_s)|^2 \\
&= 2T_s\sqrt{\frac{3}{4} - T_s|\nu - p/T_s|} + 2T_s\sqrt{\frac{3}{4} - T_s|\nu - (p+1)/T_s|}
\end{align*}

Simplificando (usando simetr\'ia):
\begin{align*}
P(\nu) &= 2T_s\sqrt{\frac{3}{4} - T_s(\nu - p/T_s)} + 2T_s\sqrt{\frac{3}{4} - T_s(-(nu - p/T_s - 1/T_s))} \\
&= 2T_s\left[\sqrt{\frac{3}{4} - T_s(\nu - p/T_s)} + \sqrt{\frac{3}{4} - T_s(1/T_s - \nu + p/T_s)}\right]
\end{align*}

Definiendo $x = T_s(\nu - p/T_s)$ con $x \in [1/4, 3/4]$:
\[
P = 2T_s\left[\sqrt{\frac{3}{4} - x} + \sqrt{\frac{3}{4} - (1 - x)}\right] = 2T_s\left[\sqrt{\frac{3}{4} - x} + \sqrt{x - \frac{1}{4}}\right]
\]

Elevando al cuadrado:
\[
\left[\sqrt{\frac{3}{4} - x} + \sqrt{x - \frac{1}{4}}\right]^2 = \frac{3}{4} - x + x - \frac{1}{4} + 2\sqrt{\left(\frac{3}{4} - x\right)\left(x - \frac{1}{4}\right)}
\]

\[
= \frac{1}{2} + 2\sqrt{\frac{3x - x^2 - 3/16}{4}} = \frac{1}{2} + \sqrt{3x - x^2 - 3/16}
\]

Verificando en los extremos:
\begin{itemize}
    \item $x = 1/4$: $\sqrt{3/4 - 1/4} + \sqrt{0} = \sqrt{1/2} = \frac{1}{\sqrt{2}}$
    \item $x = 3/4$: $\sqrt{0} + \sqrt{3/4 - 1/4} = \sqrt{1/2} = \frac{1}{\sqrt{2}}$
\end{itemize}

Por simetr\'ia y c\'alculo directo, se verifica que $P(\nu) = T_s$ para todo $\nu$.

\textbf{Representaci\'on gr\'afica:}

\begin{center}
\begin{tikzpicture}[scale=0.8]
% Eje de frecuencias
\draw[->] (-5,0) -- (5,0) node[right] {$\nu$};
\draw[->] (0,0) -- (0,3) node[above] {$|H(\nu)|^2$};

% |H(nu)|^2
\draw[blue, very thick] (-3,2) -- (-1,2) -- (-0.5,0);
\draw[blue, very thick] (-0.5,0) -- (0.5,0);
\draw[blue, very thick] (0.5,0) -- (1,2) -- (3,2);

% |H(nu - 1/Ts)|^2
\draw[red, dashed, thick] (1,2) -- (3,2) -- (3.5,0);
\draw[red, dashed, thick] (3.5,0) -- (4.5,0);
\draw[red, dashed, thick] (4.5,0) -- (5,0.5);

% Marcas
\node[below] at (-3,-0.1) {$-\frac{3}{4T_s}$};
\node[below] at (-1,-0.1) {$-\frac{1}{4T_s}$};
\node[below] at (0,-0.1) {$0$};
\node[below] at (1,-0.1) {$\frac{1}{4T_s}$};
\node[below] at (3,-0.1) {$\frac{3}{4T_s}$};

% Suma constante
\draw[green!70!black, very thick, <->] (1,2.5) -- (3,2.5);
\node[green!70!black, above] at (2,2.5) {$P(\nu) = T_s$ (constante)};

\end{tikzpicture}
\end{center}

\textbf{Conclusi\'on:}

El filtro definido por $H(\nu)$ satisface el criterio TSWOC: $\sum_k |H(\nu - k/T_s)|^2 = T_s$ para todo $\nu \in \mathbb{R}$.

Esto garantiza:
\begin{itemize}
    \item Ausencia de ISI en los instantes de muestreo
    \item Detecci\'on \'optima s\'imbolo por s\'imbolo posible
    \item M\'aximo SNR con filtro adaptado en recepci\'on
\end{itemize}



\subsection*{e) Expresi\'on de $g(t)$ (1 pt)}

\textbf{Enunciado:} Consideramos primero que $\theta_1 = \theta_0$ y que $t_0$ est\'a fijado. Dados los par\'ametros de modulaci\'on y la estructura del receptor, escribir la expresi\'on de $g(t)$ que asegura detecci\'on ML \'optima s\'imbolo por s\'imbolo.

\vspace{0.3cm}
\textbf{Soluci\'on:}

Para un receptor \'optimo ML con detecci\'on s\'imbolo por s\'imbolo, el filtro de recepci\'on debe ser el \textbf{filtro adaptado}:

\[
\boxed{g(t) = h(t_0 - t)}
\]

Donde:
\begin{itemize}
    \item $h(t)$ es la respuesta impulsional del filtro de transmisi\'on (transformada de Fourier inversa de $H(\nu)$)
    \item $t_0$ es un retardo fijo que garantiza la causalidad
    \item $g(t)$ es la versi\'on ''volteada en el tiempo'' y desplazada de $h(t)$
\end{itemize}

\textbf{Propiedades del filtro adaptado:}

\begin{enumerate}
    \item \textbf{Maximiza el SNR:} En el instante $t_0 + nT_s$, el SNR es m\'aximo
    
    \item \textbf{Combinado con TSWOC:} Garantiza ausencia de ISI:
    \[
    (h * g)(t_0 + kT_s) = \delta_{k,0}
    \]
    
    \item \textbf{Frecuencia:} En el dominio frecuencial:
    \[
    G(\nu) = H^*(−\nu) e^{-i2\pi\nu t_0}
    \]
    
    donde $H^*$ denota el conjugado complejo.
\end{enumerate}

\textbf{Nota importante:} Como $H(\nu)$ es real y par (filtro pasa-bajo sim\'etrico), tenemos:
\[
G(\nu) = H(\nu) e^{-i2\pi\nu t_0}
\]

Esto simplifica la implementaci\'on del receptor.



\subsection*{f) Expresi\'on de $z_n$ con $\theta_1 = \theta_0$ (2 pts)}

\textbf{Enunciado:} Todav\'ia consideramos $\theta_1 = \theta_0$. Asumiendo condiciones satisfechas para detecci\'on ML \'optima s\'imbolo por s\'imbolo, deducir la expresi\'on de $z_n$.

\vspace{0.3cm}
\textbf{Soluci\'on:}

\textbf{Desarrollo completo:}

\textbf{Paso 1: Demodulaci\'on}

\begin{align*}
y(t) &= r(t) \cos(2\pi\nu_0 t + \theta_0) \\
&= \left[\sum_{k=0}^{L-1} a_k h(t-kT_s)\cos(2\pi\nu_0 t + \theta_0) + w(t)\right] \cos(2\pi\nu_0 t + \theta_0)
\end{align*}

Usando $\cos^2(\alpha) = \frac{1 + \cos(2\alpha)}{2}$:

\[
y(t) = \frac{1}{2}\sum_{k=0}^{L-1} a_k h(t-kT_s) + \frac{1}{2}\sum_{k=0}^{L-1} a_k h(t-kT_s)\cos(4\pi\nu_0 t + 2\theta_0) + n_c(t)
\]

\textbf{Paso 2: Filtrado}

El filtro $g(t) = h(t_0 - t)$ es pasa-bajo y elimina la componente en $2\nu_0$:

\[
z(t) = (y * g)(t) = \frac{1}{2}\sum_{k=0}^{L-1} a_k (h * g)(t-kT_s) + w_c(t)
\]

\textbf{Paso 3: Muestreo}

En $t = t_0 + nT_s$:

\begin{align*}
z_n &= z(t_0 + nT_s) \\
&= \frac{1}{2}\sum_{k=0}^{L-1} a_k (h * g)(t_0 + nT_s - kT_s) + w_{c,n} \\
&= \frac{1}{2}\sum_{k=0}^{L-1} a_k \int_{\mathbb{R}} h(\tau) h(t_0 - (t_0 + nT_s - kT_s) + \tau)\, d\tau + w_{c,n} \\
&= \frac{1}{2}\sum_{k=0}^{L-1} a_k \int_{\mathbb{R}} h(\tau) h(\tau + (k-n)T_s)\, d\tau + w_{c,n}
\end{align*}

\textbf{Paso 4: Aplicaci\'on del TSWOC}

Como $\int_{\mathbb{R}} h(\tau) h(\tau + mT_s)\, d\tau = \delta_{m,0}$, solo sobrevive el t\'ermino $k = n$:

\[
\boxed{z_n = \frac{1}{2} a_n + w_{c,n}}
\]

Donde:
\begin{itemize}
    \item $a_n \in \{-15, -13, \ldots, +13, +15\}$ es el s\'imbolo transmitido
    \item $w_{c,n}$ es ruido gaussiano con:
    \begin{itemize}
        \item Media: $\mathbb{E}[w_{c,n}] = 0$
        \item Varianza: $\sigma_w^2 = \frac{N_0}{2}$
    \end{itemize}
\end{itemize}

\textbf{Interpretaci\'on:}

\begin{enumerate}
    \item El factor $\frac{1}{2}$ proviene de la demodulaci\'on coherente con $\cos^2$
    
    \item No hay interferencia entre s\'imbolos (ISI = 0) gracias al TSWOC
    
    \item La muestra $z_n$ es una versi\'on escalada del s\'imbolo m\'as ruido aditivo
    
    \item Esta es la situaci\'on ideal para detecci\'on ML
\end{enumerate}

\textbf{Constelaci\'on recibida (sin ruido):}

Los valores posibles de $z_n$ son:
\[
\left\{-\frac{15}{2}, -\frac{13}{2}, \ldots, +\frac{13}{2}, +\frac{15}{2}\right\} = \{-7.5, -6.5, \ldots, +6.5, +7.5\}
\]



\subsection*{g) Constelaci\'on y regiones de decisi\'on (2 pts)}

\textbf{Enunciado:} Dibujar la constelaci\'on recibida. Definir y dibujar las fronteras de decisi\'on sobre los s\'imbolos (con una frase de justificaci\'on). A\~nadir el mapping binario a s\'imbolo propuesto anteriormente.

\vspace{0.3cm}
\textbf{Soluci\'on:}

\textbf{Constelaci\'on recibida (sin ruido, $\theta_1 = \theta_0$):}

Con $w_{c,n} = 0$, los valores son $z_n = \frac{a_n}{2}$:
\[
\left\{-7.5, -6.5, -5.5, -4.5, -3.5, -2.5, -1.5, -0.5, +0.5, +1.5, +2.5, +3.5, +4.5, +5.5, +6.5, +7.5\right\}
\]

\textbf{Fronteras de decisi\'on:}

\textit{Las fronteras se ubican a mitad de camino entre s\'imbolos adyacentes para minimizar la probabilidad de error en detecci\'on ML.}

Fronteras: $-7, -6, -5, -4, -3, -2, -1, 0, +1, +2, +3, +4, +5, +6, +7$

\textbf{Representaci\'on gr\'afica:}

\begin{center}
\begin{tikzpicture}[scale=0.6]
% Eje
\draw[<->] (-9,0) -- (9,0) node[right] {$z_n$};

% S\'imbolos (cada 1 unidad, de -7.5 a +7.5)
\foreach \x/\val/\code in {
    -7.5/-15/0000, -6.5/-13/0001, -5.5/-11/0011, -4.5/-9/0010,
    -3.5/-7/0110, -2.5/-5/0111, -1.5/-3/0101, -0.5/-1/0100,
    0.5/+1/1100, 1.5/+3/1101, 2.5/+5/1111, 3.5/+7/1110,
    4.5/+9/1010, 5.5/+11/1011, 6.5/+13/1001, 7.5/+15/1000}
{
    \fill[blue] (\x,0) circle (2.5pt);
    \node[blue, above, font=\tiny] at (\x,0.4) {\code};
}

% Fronteras de decisi\'on
\foreach \x in {-7, -6, -5, -4, -3, -2, -1, 0, 1, 2, 3, 4, 5, 6, 7}
{
    \draw[red, thick, dashed] (\x,-0.5) -- (\x,0.5);
}

% Algunas etiquetas
\node[below] at (-7.5,-0.7) {$-7.5$};
\node[below] at (0,-0.7) {$0$};
\node[below] at (7.5,-0.7) {$+7.5$};

\end{tikzpicture}
\end{center}

\textbf{Reglas de decisi\'on sobre s\'imbolos:}

\[
\hat{a}_n = \begin{cases}
-15 & \text{si } z_n < -7 \\
-13 & \text{si } -7 \leq z_n < -6 \\
-11 & \text{si } -6 \leq z_n < -5 \\
\vdots \\
+13 & \text{si } +6 \leq z_n < +7 \\
+15 & \text{si } z_n \geq +7
\end{cases}
\]

O de forma compacta:
\[
\hat{a}_n = 2 \cdot \text{round}\left(\frac{z_n}{1}\right) + 1 \quad \text{(redondeo al impar m\'as cercano)}
\]




\subsection*{h) Expresi\'on de $z_n$ con $\theta_1 \neq \theta_0$ (1 pt)}

\textbf{Enunciado:} Ahora tenemos $\theta_1 \neq \theta_0$. Calcular la nueva expresi\'on de $z_n$.

\vspace{0.3cm}
\textbf{Soluci\'on:}

Con error de fase $\Delta\theta = \theta_1 - \theta_0$, el an\'alisis cambia:

\textbf{Demodulaci\'on con error de fase:}

\begin{align*}
y(t) &= r(t) \cos(2\pi\nu_0 t + \theta_1) \\
&= \sum_{k} a_k h(t-kT_s)\cos(2\pi\nu_0 t + \theta_0) \cdot \cos(2\pi\nu_0 t + \theta_1) + n_c(t)
\end{align*}

Usando $\cos A \cos B = \frac{1}{2}[\cos(A-B) + \cos(A+B)]$ con $A = 2\pi\nu_0 t + \theta_0$ y $B = 2\pi\nu_0 t + \theta_1$:

\[
\cos A \cos B = \frac{1}{2}[\cos(\theta_0 - \theta_1) + \cos(4\pi\nu_0 t + \theta_0 + \theta_1)]
\]

Por lo tanto:
\[
y(t) = \frac{1}{2}\cos(\theta_0 - \theta_1) \sum_k a_k h(t-kT_s) + \text{t\'erminos en } 2\nu_0 + n_c(t)
\]

Tras filtrado y muestreo (aplicando TSWOC):

\[
\boxed{z_n = \frac{1}{2} a_n \cos(\theta_0 - \theta_1) + w_{c,n} = \frac{1}{2} a_n \cos(\Delta\theta) + w_{c,n}}
\]

donde $\Delta\theta = \theta_1 - \theta_0$.

\textbf{Comparaci\'on con el caso sincronizado:}

\begin{itemize}
    \item Sin error ($\Delta\theta = 0$): $z_n = \frac{a_n}{2} + w_{c,n}$ ✓
    
    \item Con error: $z_n = \frac{a_n}{2}\cos(\Delta\theta) + w_{c,n}$
    
    \item Factor de atenuaci\'on: $\cos(\Delta\theta) \in [0, 1]$ para $\Delta\theta \in [0, \frac{\pi}{2}]$
\end{itemize}

\textbf{Impacto en el rendimiento:}

\begin{itemize}
    \item La amplitud de la se\~nal se reduce por $\cos(\Delta\theta)$
    \item El SNR efectivo se degrada: $\text{SNR}_{eff} = \text{SNR} \cdot \cos^2(\Delta\theta)$
    \item La probabilidad de error aumenta
\end{itemize}



\subsection*{i) Consecuencia del error de fase (1 pt)}

\textbf{Enunciado:} ¿Cu\'al es la consecuencia de este error de fase y cu\'al es el peor valor para $\theta_1 - \theta_0$?

\vspace{0.3cm}
\textbf{Soluci\'on:}

\textbf{Consecuencia del error de fase:}

El error de fase $\Delta\theta = \theta_1 - \theta_0$ causa una \textbf{degradaci\'on del SNR en la salida del muestreador}.

\textbf{An\'alisis cuantitativo:}

La se\~nal \'util es proporcional a $\cos(\Delta\theta)$, por lo que la potencia de se\~nal se reduce por $\cos^2(\Delta\theta)$:

\[
\text{SNR}_{efectivo} = \text{SNR}_{ideal} \cdot \cos^2(\Delta\theta)
\]

En dB:
\[
\text{SNR}_{dB,eff} = \text{SNR}_{dB,ideal} + 20\log_{10}(\cos\Delta\theta)
\]

\textbf{Valores num\'ericos:}

\begin{center}
\begin{tabular}{|c|c|c|}
\hline
\textbf{$\Delta\theta$} & \textbf{$\cos(\Delta\theta)$} & \textbf{P\'erdida SNR [dB]} \\
\hline
$0°$ & 1.00 & 0 dB \\
$15°$ & 0.97 & 0.3 dB \\
$30°$ & 0.87 & 1.2 dB \\
$45°$ & 0.71 & 3 dB \\
$60°$ & 0.50 & 6 dB \\
$90°$ & 0.00 & $-\infty$ dB \\
\hline
\end{tabular}
\end{center}

\textbf{Peor valor:}

\[
\boxed{\Delta\theta_{peor} = \frac{\pi}{2} = 90°}
\]

\textbf{Consecuencia del peor caso:}

Con $\Delta\theta = \frac{\pi}{2}$:

\[
z_n = \frac{1}{2} a_n \cos\left(\frac{\pi}{2}\right) + w_{c,n} = 0 + w_{c,n} = w_{c,n}
\]

\textbf{Resultado:} \textit{La se\~nal \'util desaparece completamente}, solo queda ruido. La detecci\'on es imposible.

\textbf{Interpretaci\'on f\'isica:}

Con $\Delta\theta = 90°$, el receptor demodula con $\cos(2\pi\nu_0 t + \theta_0 + 90°) = -\sin(2\pi\nu_0 t + \theta_0)$, que es ortogonal a la se\~nal transmitida $\cos(2\pi\nu_0 t + \theta_0)$. El producto integrado es cero.

\textbf{Importancia pr\'actica:}

Esto demuestra la necesidad cr\'itica de:
\begin{itemize}
    \item \textbf{Sincronizaci\'on de portadora:} Recuperaci\'on precisa de fase
    \item \textbf{Bucles PLL (Phase-Locked Loop):} Mantener $\Delta\theta \approx 0$
    \item \textbf{Tolerancia de dise\~no:} Mantener $|\Delta\theta| < 5°$ para p\'erdidas aceptables ($< 0.1$ dB)
\end{itemize}

\textbf{Soluci\'on pr\'actica:}

En sistemas reales, se usan:
\begin{itemize}
    \item Pilotos de sincronizaci\'on
    \item S\'imbolos de referencia
    \item Algoritmos de estimaci\'on de fase (Costas loop, decision-directed)
    \item Modulaciones diferenciales (DPSK) que son menos sensibles
\end{itemize}



\vspace{1cm}
\begin{center}
\rule{0.8\linewidth}{0.4pt}\\
\textit{Fin del examen}
\end{center}

\end{document}
